\section{Ablations (claim D1)}
\label{sec:ablations}

Each ablation removes one component, and its forecasts are compared with
the full model on the same origins. The run covers all
\R{ho/abldata/accounts_with_eligible_origin} held-out accounts with eligible
origins (Table~\ref{tab:ablations}). Two caveats come first.
\begin{enumerate}[label=(\roman*)]
\item \textbf{The components were removed from V0, not from V2.} D1 was
  written into the plan before V2 existed. The ablations therefore show what
  each component does inside the original model. They do not show what it
  does inside V2, whose spending model replaces the smooth rate and the
  lumps.
\item \textbf{The alarm-defined strata are circular} (\S\ref{sec:regime}).
  They are selected by the full model's own detector.
\end{enumerate}
They were also run after the main held-out results were known, in a
second logged run, because the first run omitted them
(\S\ref{sec:timeline}).

\begin{table}[t]
\centering\scriptsize\setlength{\tabcolsep}{3.5pt}
\tab{heldout-ablations}
\caption{Held-out ablations against V0. Skill: the ablation's CRPS skill
against V0 (positive: the ablation is better). $\Delta$cov90: the
ablation's \lC\% coverage minus V0's, in points. Columns are horizons
(days). Intervals are in the result files.}
\label{tab:ablations}
\end{table}

\begin{itemize}
\item \textbf{Adaptive noise off (fixed priors).} The prediction was worse
  coverage. Coverage does fall at \hA\ and \hC\ days
  (\R{ho/ablp/all/abl-fixed-noise/7/cov90_diff} and
  \R{ho/ablp/all/abl-fixed-noise/90/cov90_diff} points) and is unchanged
  at \hB, so the prediction holds. But CRPS is \emph{better} without
  adaptation: skill \R{ho/ablp/all/abl-fixed-noise/30/skill@3} and
  \R{ho/ablp/all/abl-fixed-noise/90/skill@3} at \hB\ and \hC\ days.
\item \textbf{Lumps off.} The prediction was under-coverage at the wide
  levels, mostly at long horizons. It holds strongly: \lC\% coverage falls
  by \R{ho/ablp/all/abl-no-lumps/30/cov90_diff} and
  \R{ho/ablp/all/abl-no-lumps/90/cov90_diff} points at \hB\ and \hC\ days.
  On sparse Berka spending, the lump resampling carries much of V0's
  spread.
\item \textbf{Spend floor off.} The prediction was no change except where
  the floor binds. It holds overall: effects stay below one point of
  coverage and below one percent of CRPS. After alarms, CRPS is worse
  without the floor
  (\R{ho/ablp/regime-alarm-in-last-28-days/abl-no-floor/90/skill@3} at
  \hC\ days), which fits the floor binding where the estimated rate has
  just jumped.
\item \textbf{Detector off.} Both predicted signs are reversed
  (\S\ref{sec:regime}). By the plan's rule as written, this falsifies
  the detector's claim. Because both strata are selected by the detector's
  own alarms, we read the result as descriptive, not as evidence about
  true changes.
\item \textbf{Noise-variance draws off.} No prediction was registered.
  The effects are negligible:
  \R{ho/ablp/all/abl-no-var-draw/90/skill@3} in skill and
  \R{ho/ablp/all/abl-no-var-draw/90/cov90_diff} points of coverage at
  \hC\ days.
\end{itemize}
